\documentclass[12pt,a4paper]{article}
\usepackage[UTF8]{ctex}
\usepackage{amsmath,amssymb,amsthm}
\usepackage{geometry}

\geometry{margin=2.5cm}

\title{为什么P控制器存在稳态误差？}
\author{第11章补充说明}
\date{}

\begin{document}

\maketitle

\section{问题}

\textbf{为什么P控制器（比例控制器）在时间$t \to \infty$时，误差$e$不为0？}

\section{P控制器回顾}

P控制器的控制律为：
\begin{equation}
\tau = K_p(\theta_d - \theta) = K_p \theta_e
\end{equation}
其中：
\begin{itemize}
\item $\tau$是控制力矩
\item $K_p > 0$是比例增益
\item $\theta_d$是期望位置
\item $\theta$是实际位置
\item $\theta_e = \theta_d - \theta$是位置误差
\end{itemize}

\section{稳态误差的根本原因}

\subsection{关键洞察}

\textbf{P控制器存在稳态误差的根本原因：}

当系统达到稳态时，如果存在\textbf{持续的外部干扰}（如重力、摩擦力、负载等），P控制器需要\textbf{非零误差}来产生\textbf{非零的控制输出}以抵消这些干扰。

如果误差为0，控制输出也为0，无法抵消干扰，系统无法保持平衡。

\subsection{数学分析}

考虑一个简单的单关节机器人系统，动力学方程为：
\begin{equation}
M\ddot{\theta} = \tau - mgr\cos\theta - b\dot{\theta}
\end{equation}
其中：
\begin{itemize}
\item $M$是惯性
\item $\tau$是控制力矩
\item $mgr\cos\theta$是重力项（持续干扰）
\item $b\dot{\theta}$是摩擦力项
\end{itemize}

使用P控制器：
\begin{equation}
\tau = K_p(\theta_d - \theta) = K_p \theta_e
\end{equation}

\subsubsection{稳态分析}

在稳态时（$t \to \infty$）：
\begin{itemize}
\item $\ddot{\theta} = 0$（加速度为零）
\item $\dot{\theta} = 0$（速度为零，假设没有持续运动）
\item $\theta = \theta_{ss}$（稳态位置）
\item $\theta_e = \theta_{e,ss}$（稳态误差）
\end{itemize}

稳态时的动力学方程变为：
\begin{equation}
0 = \tau_{ss} - mgr\cos\theta_{ss} - 0
\end{equation}

即：
\begin{equation}
\tau_{ss} = mgr\cos\theta_{ss}
\end{equation}

但控制力矩为：
\begin{equation}
\tau_{ss} = K_p \theta_{e,ss}
\end{equation}

因此：
\begin{equation}
K_p \theta_{e,ss} = mgr\cos\theta_{ss}
\end{equation}

如果$\theta_{ss} \approx \theta_d$（稳态位置接近期望位置），则：
\begin{equation}
\theta_{e,ss} = \frac{mgr\cos\theta_d}{K_p}
\end{equation}

\textbf{结论：}稳态误差$\theta_{e,ss} \neq 0$，除非$mgr\cos\theta_d = 0$（无重力或重力被完全补偿）。

\section{详细分析：不同情况}

\subsection{情况1：无干扰系统}

如果系统没有持续干扰（如重力、摩擦力等），P控制器可以达到零稳态误差。

\textbf{动力学：}
\begin{equation}
M\ddot{\theta} = \tau
\end{equation}

\textbf{控制律：}
\begin{equation}
\tau = K_p(\theta_d - \theta)
\end{equation}

\textbf{误差动力学：}
\begin{equation}
M\ddot{\theta}_e = -K_p \theta_e
\end{equation}

这是一个二阶系统，如果$K_p > 0$，误差会收敛到0。

\textbf{结论：}无干扰时，P控制器可以达到零稳态误差。

\subsection{情况2：有持续干扰（重力）}

\textbf{动力学：}
\begin{equation}
M\ddot{\theta} = \tau - mgr\cos\theta
\end{equation}

\textbf{控制律：}
\begin{equation}
\tau = K_p(\theta_d - \theta)
\end{equation}

\textbf{稳态时：}
\begin{align}
0 &= K_p(\theta_d - \theta_{ss}) - mgr\cos\theta_{ss} \\
K_p(\theta_d - \theta_{ss}) &= mgr\cos\theta_{ss} \\
K_p \theta_{e,ss} &= mgr\cos\theta_{ss}
\end{align}

如果$\theta_{ss} \approx \theta_d$，则：
\begin{equation}
\theta_{e,ss} \approx \frac{mgr\cos\theta_d}{K_p}
\end{equation}

\textbf{结论：}有持续干扰时，P控制器存在稳态误差，误差大小与干扰成正比，与增益$K_p$成反比。

\subsection{情况3：有摩擦力}

\textbf{动力学：}
\begin{equation}
M\ddot{\theta} = \tau - b\dot{\theta}
\end{equation}

\textbf{控制律：}
\begin{equation}
\tau = K_p(\theta_d - \theta)
\end{equation}

\textbf{稳态时：}
\begin{align}
0 &= K_p(\theta_d - \theta_{ss}) - b \cdot 0 \\
K_p(\theta_d - \theta_{ss}) &= 0 \\
\theta_{e,ss} &= 0
\end{align}

\textbf{结论：}如果只有速度相关的摩擦力（无静摩擦），稳态时速度为零，摩擦力也为零，P控制器可以达到零稳态误差。

但如果存在\textbf{静摩擦}（Coulomb摩擦）：
\begin{equation}
M\ddot{\theta} = \tau - F_c \text{sign}(\dot{\theta})
\end{equation}

在稳态时，如果$\dot{\theta} = 0$，静摩擦可能不为零，需要非零控制力矩来抵消，因此存在稳态误差。

\section{物理直观理解}

\subsection{类比：弹簧-质量系统}

想象一个弹簧-质量系统，受到持续的重力作用：

\begin{itemize}
\item 弹簧力：$F_{\text{spring}} = kx$（$x$是位移，$k$是弹簧常数）
\item 重力：$F_{\text{gravity}} = mg$（持续向下）
\item 平衡时：$kx = mg$，因此$x = mg/k \neq 0$
\end{itemize}

P控制器就像弹簧：控制力矩$\tau = K_p \theta_e$就像弹簧力$F = kx$。

如果存在持续干扰（重力），系统会在非零误差处达到平衡，误差产生的控制力矩正好抵消干扰。

\subsection{能量平衡观点}

在稳态时：
\begin{itemize}
\item 控制力矩做的功 = 干扰力矩做的功
\item 如果误差为0，控制力矩为0，无法抵消干扰
\item 因此，必须存在非零误差来产生非零控制力矩
\end{itemize}

\section{稳态误差的表达式}

对于一般系统，稳态误差可以表示为：

\begin{equation}
e_{ss} = \lim_{t \to \infty} e(t) = \frac{\text{持续干扰}}{K_p}
\end{equation}

更精确地，对于单位阶跃输入和持续干扰$d$：

\begin{equation}
e_{ss} = \frac{d}{K_p}
\end{equation}

\textbf{关键观察：}
\begin{itemize}
\item 稳态误差与持续干扰成正比
\item 稳态误差与增益$K_p$成反比
\item 增大$K_p$可以减小稳态误差，但不能完全消除（除非$K_p \to \infty$，这不现实）
\end{itemize}

\section{为什么增大Kp不能完全消除稳态误差？}

理论上，如果$K_p \to \infty$，则$e_{ss} \to 0$。但在实际中：

\begin{enumerate}
\item \textbf{物理限制：}执行器有最大力矩限制，$K_p$不能无限大
\item \textbf{稳定性问题：}大$K_p$可能导致系统不稳定（如离散时间系统中的问题）
\item \textbf{噪声放大：}大$K_p$会放大传感器噪声
\item \textbf{速度限制：}大$K_p$可能产生超出物理限制的速度指令
\end{enumerate}

因此，在实际应用中，$K_p$不能无限增大，稳态误差无法完全消除。

\section{解决方案：使用PI控制器}

为了消除稳态误差，需要使用\textbf{积分项}（I项）。

\subsection{PI控制器}

PI控制器的控制律为：
\begin{equation}
\tau = K_p \theta_e + K_i \int_0^t \theta_e(\tau) d\tau
\end{equation}

\textbf{为什么积分项可以消除稳态误差？}

\begin{itemize}
\item 即使误差很小，积分项会持续累积
\item 只要误差不为零，积分项就会增长
\item 积分项产生的控制力矩可以抵消持续干扰
\item 当误差为零时，积分项保持常数，提供抵消干扰所需的力矩
\end{itemize}

\textbf{稳态分析：}

在稳态时，如果$\theta_e = 0$，则：
\begin{equation}
\tau_{ss} = K_p \cdot 0 + K_i \int_0^{\infty} \theta_e(\tau) d\tau = K_i \cdot \text{常数}
\end{equation}

这个常数可以正好抵消持续干扰，使得$\theta_e = 0$成为可能的稳态。

\section{数值示例}

\subsection{示例1：有重力的系统}

设：
\begin{itemize}
\item $mgr = 10$ N·m（重力力矩）
\item $K_p = 100$ N·m/rad
\item $\theta_d = 0$ rad
\end{itemize}

稳态误差：
\begin{equation}
\theta_{e,ss} = \frac{10}{100} = 0.1 \text{ rad} \approx 5.7^\circ
\end{equation}

如果增大$K_p = 1000$：
\begin{equation}
\theta_{e,ss} = \frac{10}{1000} = 0.01 \text{ rad} \approx 0.57^\circ
\end{equation}

误差减小了，但仍然存在。

\subsection{示例2：使用PI控制器}

使用PI控制器：
\begin{equation}
\tau = 100 \theta_e + 10 \int_0^t \theta_e(\tau) d\tau
\end{equation}

在稳态时，如果$\theta_e = 0$：
\begin{equation}
\tau_{ss} = 10 \int_0^{\infty} \theta_e(\tau) d\tau
\end{equation}

积分项会调整到$\tau_{ss} = 10$ N·m，正好抵消重力，使得$\theta_e = 0$。

\section{总结}

\textbf{为什么P控制器存在稳态误差？}

\begin{enumerate}
\item \textbf{根本原因：}当存在持续外部干扰（如重力、静摩擦）时，P控制器需要非零误差来产生非零控制力矩以抵消干扰。

\item \textbf{数学表达：}稳态误差$e_{ss} = \frac{\text{持续干扰}}{K_p}$

\item \textbf{物理直观：}就像弹簧-质量系统，在重力作用下，弹簧必须伸长（非零位移）才能产生抵消重力的力。

\item \textbf{解决方案：}使用PI控制器（加入积分项）可以消除稳态误差，因为积分项可以在误差为零时仍提供抵消干扰所需的控制力矩。
\end{enumerate}

\textbf{关键点：}
\begin{itemize}
\item P控制器在无干扰系统中可以达到零稳态误差
\item P控制器在有持续干扰的系统中必然存在稳态误差
\item 增大$K_p$可以减小稳态误差，但不能完全消除（受物理限制）
\item PI控制器可以消除稳态误差
\end{itemize}

\end{document}

